\section{Introduction}
\label{sec:intro}

% Retrieval-Augmented Generation (RAG) has become a standard paradigm for augmenting the capabilities of large language models (LLMs), particularly n settings involving with out-of-distribution or domain-specific corpora~\cite{rag_2020,rag_survey}. Despite strong performance on generic benchmarks, these evaluations often fail to reflect the specific needs of domain experts across disciplines~\cite{Chen_Lin_Han_Sun_2024,domainrag}.

% This limitation is critical in complex domains such as historical research, where question answering involves complex reasoning patterns, including multi-hop inference across heterogeneous sources, imbalance corpora~\cite{historiqa-pellet}.

% Historians do not pose interchangeable questions. They have specific needs, intentions, and thematic concerns that shape how they look for evidence. This diversity justifies task-specific approaches to RAG, whether in the indexing, retrieval, or answer generation phase.

Retrieval-Augmented Generation (RAG) has become a standard paradigm for
augmenting large language models (LLMs), particularly in settings
involving out-of-distribution or domain-specific
corpora~\cite{rag_2020,rag_survey}. Despite strong performance on
generic benchmarks, these evaluations often fail to reflect the
specific needs of domain experts~\cite{Chen_Lin_Han_Sun_2024,domainrag},
a limitation that becomes critical in complex settings such as
historical research, where question answering requires multi-hop
inference across heterogeneous, imbalanced
sources~\cite{historiqa-pellet}. Historians do not pose interchangeable
questions: their specific intentions and thematic concerns shape how
they look for evidence, which justifies task-specific approaches to
RAG across the indexing, retrieval, and generation phases.


%Researchers such as historians do not pose interchangeable questions. They have specific needs, intentions, and thematic concerns that shape how they look for evidence. This diversity justifies task-specific approaches to RAG, whether in the indexing, retrieval, or answer generation phase.

% Standard RAG pipelines, even when enhanced with reranking, exhibit significant limitations in such settings. These limitations are exacerbated by highly imbalanced corpora, where retrieval systems tend to favour dominant sources, leading to biased or incomplete evidence aggregation. 

 % Recent work~\cite{hipporagv2,adaptive_chunking} has begun to address some of these issues by revisiting indexing and document segmentation strategies. However, two weaknesses persist: chunking and retrieval are addressed independently of each other, and both are treated independently of the query. This implicitly assumes that a single segmentation strategy and a unified retrieval scheme can perform well across all query types -- an assumption that, as we show, does not hold on complex historical corpora.

 % Our proposition rests on two intertwined claims. First, the optimisation of chunking and retrieval in RAG admits \emph{no one-size-fits-all solution}. Despite having an overall best chunking strategy, adding a query based chunk optimisation yields significativly better results. More generally the best segmentation, the best metadata filter, and the best source selection vary systematically with the type of question being asked.
 % Therefore, and it's our second claim chunking and retrieval strategies can and must be partly conditioned online on the user query, while the heavy lifting (index construction, strategy discovery) remains offline and conditioned on the document. Concretely, we precompute a small set of indices offline, and at inference time we route each query, at negligible online cost, to the index and retrieval configuration best suited to it. This combination of rejecting one-size-fits-all and splitting the optimisation across an offline document-conditioned axis and an online query-conditioned axis is the core of our proposal.

Recent work~\cite{hipporagv2,adaptive_chunking} has begun to revisit
indexing and document segmentation strategies. However, two weaknesses
persist: chunking and retrieval are addressed independently of each
other, and both are treated independently of the query. This
\emph{one-size-fits-all} assumption that a single segmentation
strategy and a unified retrieval scheme perform well across all query
types does not hold on complex historical corpora. We therefore
argue that the optimisation should be split along two axes: an
\emph{offline}, document-conditioned axis (index construction, strategy
discovery) and an \emph{online}, query-conditioned axis that routes
each query, at negligible cost, to the index and retrieval
configuration best suited to it.


% We operationalize this position with three contributions, the first two acting on the index side (offline strategy learning, online routing) and the third on the retrieval side (online source selection):

% \begin{itemize}


%   \item \textbf{Query-conditioned chunking.}
%     We introduce a framework that discovers, offline, semantic clusters of questions and learns the segmentation strategy that maximises retrieval coverage for each cluster. At inference, on a held out test set, each query is projected into the learned embedding space and routed via nearest-centroid assignment to the pre-built index corresponding to its cluster. The marginal online cost is one projection and one nearest-neighbour lookup.
%     This yields this realises +7.5 pp Coverage@3 over
% the global \texttt{10000_hierarchical} baseline.
 
%   \item \textbf{Query-conditioned metadata strategies.}
%      Reusing the same cluster topology, we learn per cluster the
% optimal metadata configuration (filter / rerank / hybrid) on \textit{Les Débats}. This adds +10.8 pp
% Coverage@3 under a Naive Retriever and remains positive (+1.0 to +2.5 pp) under sophisticated
% retrievers.

%   \item \textbf{Multi-source rounting (QRe + UMS).}
%      A supervised classifier predicts which corpora are required for
% the query — empirically a binary press-vs-parliamentary decision (99.4\% accuracy). Coupled with a
% uniform per-source allocation (UMS), the combined retriever reaches Recall@3 of 52.3 (+16.4 pp over a
% Naive Retriever, +20.6 pp over HippoRAGv2). We show that QRe and UMS attack structurally distinct
% failure modes (source contamination vs.\ source starvation) and compose strictly
% \end{itemize}

% Taken together, these three mechanisms outperform both a naive RAG baseline and state-of-the-art methods that rely on a generalized, task-agnostic approach. Full implementation details are available in the code repository.\footnote{Code available at \href{https://github.com/atomegoyan/debattre-generate-questions.git}{GitHub}.}

We operationalise this position with three contributions, the first
two on the index side (offline strategy learning, online routing) and
the third on the retrieval side (online source selection):

\paragraph{Query-conditioned chunking.} We discover semantic
question clusters offline, learn per cluster the segmentation
that maximises retrieval coverage, and at inference route each
query via nearest-centroid assignment to the corresponding
pre-built index ($+7.5$~pp Coverage@3 over the global
\texttt{10000\_hierarchical} baseline). Marginal online cost is
one projection and one nearest-neighbour lookup.

\paragraph{Query-conditioned metadata selection.} Reusing the
same cluster topology, we learn per cluster the optimal metadata
configuration (filter / rerank / hybrid) on \textit{Les D\'ebats}
($+10.8$~pp Coverage@3 under a Naive Retriever; $+1.0$ to
$+2.5$~pp under sophisticated retrievers).

\paragraph{Multi-source routing (QRe + UMS).} A supervised
classifier predicts which corpora are required empirically a
binary press-vs-parliamentary decision (99.4\% accuracy),
which combined with a uniform per-source allocator (UMS) reaches
Recall@3 of $52.3$ ($+16.4$ pp over the Naive Retriever,
$+20.6$ pp over HippoRAGv2). The two components attack
structurally distinct failure modes (source contamination
vs.\ source starvation) and compose strictly.

These three mechanisms outperform both a naive RAG
baseline and state-of-the-art methods that rely on a task-agnostic
approach.